(1) 电子方程 $m\dot{\mathbf v}=-e(\mathbf E+\mathbf v\times B_0\mathbf e_z)$。定义 $v_\pm=v_x\pm iv_y$、$E_\pm=E_x\pm iE_y$，因为 $(\mathbf v\times\mathbf B)_\pm=\mp iB_0v_\pm$，方程变成 $-i(\omega\pm\Omega)v_\pm=-eE_\pm/m$，故 $v_\pm=-ieE_\pm/[m(\omega\pm\Omega)]$。

(2) 电流 $j_\pm=-nev_\pm=i\varepsilon_0\omega_p^2E_\pm/(\omega\pm\Omega)$，其中 $\omega_p^2=ne^2/(\varepsilon_0m)$。结合Maxwell横向波动方程得 $\boxed{n_\pm^2=c^2k^2/\omega^2=1-\omega_p^2/[\omega(\omega\pm\Omega)]}$。$E_-$ 分支在 $\omega=\Omega$ 处有理想无碰撞回旋共振，具体左右旋命名依观察方向而定。

(3) 对共振分支 $n_-^2=1-\omega_p^2/[\omega(\omega-\Omega)]$，在 $0<\omega\ll\Omega$ 时 $n_-^2\simeq1+\omega_p^2/(\Omega\omega)$。若第二项还远大于1，便有 $\boxed{\omega\simeq\Omega c^2k^2/\omega_p^2}$，群速度 $v_g=d\omega/dk\simeq2\Omega c^2k/\omega_p^2$，是色散明显的哨声波而非恒速电磁波。

(4) 当 $\Omega\to0$，两圆偏振简并为 $n^2=1-\omega_p^2/\omega^2$，即无磁等离子体的截止色散。离子固定条件要求观察频率相对离子动力学时间尺度不太低，否则必须加入离子电流；有限碰撞会把分母改成复频移，产生正的吸收并抹平无损共振奇点。
